%%%%%%%%%%%%%%%%%%%%%%%%%%%%%%%%%%%%%%%%%%%%%%%%%%%%%%%%%%%%%%%%%%%%%%
%%  MoodTown（Agentopia 风格稿 V2）中文对照检查版
%%
%%  本文件为 agentopia-style-manuscript.tex 的中文对照检查版。
%%  生成日期：2026-09-27。
%%  用途：供作者逐段审读译文与论证结构；所有占位符以红色 \ph{} 显示。
%%  投稿请使用英文版 agentopia-style-manuscript.tex。
%%  编译：latexmk -xelatex agentopia-style-manuscript-zh.tex
%%  （需 ctex 宏包；参考文献使用同目录 agentopia-references.bib）
%%%%%%%%%%%%%%%%%%%%%%%%%%%%%%%%%%%%%%%%%%%%%%%%%%%%%%%%%%%%%%%%%%%%%%

\documentclass[11pt,a4paper]{article}

\usepackage[UTF8,heading=false,fontset=windows]{ctex}
\usepackage{geometry}
\geometry{margin=2.5cm}
\usepackage{amsmath,amssymb}
\usepackage{booktabs}
\usepackage{graphicx}
\usepackage{xcolor}
\usepackage[round]{natbib}
\usepackage{enumitem}
\usepackage[font=small,labelfont=bf]{caption}
\usepackage{titlesec}
\usepackage[colorlinks=true,linkcolor=blue!60!black,citecolor=blue!60!black,urlcolor=blue!60!black]{hyperref}

% 行距与段间距
\linespread{1.35}
\setlength{\parskip}{0.4em}
\setlength{\parindent}{2em}

% 蓝色章节标题
\titleformat{\section}{\Large\bfseries\color{blue!40!black}}{\thesection}{1em}{}
\titleformat{\subsection}{\large\bfseries\color{blue!30!black}}{\thesubsection}{1em}{}
% H4：粗体行内段首标题
\titleformat{\paragraph}[runin]{\bfseries}{}{0pt}{}
\titlespacing{\paragraph}{0pt}{0.8em}{0.6em}

% 占位符高亮（红色），便于检查时定位
\newcommand{\ph}[1]{\textcolor{red}{[#1]}}

% 平台专名（保留英文，不译）
\newcommand{\sys}{MoodTown}

% 中英文名称
\renewcommand{\abstractname}{摘要}
\renewcommand{\refname}{参考文献}
\renewcommand{\figurename}{图}
\renewcommand{\tablename}{表}

\title{\textbf{\sys：以生成式智能体模拟抑郁症社区\\并进行 \emph{in silico} 多臂 CBT 试验}}
\author{%
  \ph{作者 1}$^{1}$,\; \ph{作者 2}$^{1}$,\; \ph{作者 3}$^{1}$,\; \ph{作者 4}$^{1}$\thanks{通讯作者：\ph{email@example.org}}\\
  $^{1}$\ph{院系}，\ph{机构}，\ph{城市}，\ph{国家}
}
\date{预印本，2026 年 9 月。\quad（中文对照检查版，生成日期 2026-09-27）}

\begin{document}

\maketitle

\begin{center}
\rule{\textwidth}{0.5pt}\\[0.3em]
\small\textit{英文原题：MoodTown: Simulating a Depression Community with Generative Agents\\
and In-Silico Multi-Arm CBT Trials}\\[0.3em]
\rule{\textwidth}{0.5pt}
\end{center}

\vspace{0.5em}

%%====================================================================
%%  摘要
%%====================================================================
\begin{abstract}
\noindent
抑郁症是全球致残的主要原因之一，估计影响 2.8 亿人~\citep{who2022}。
人类在社会生活之中发展、维持抑郁并从中恢复，这引出一个自然的问题：由 LLM 驱动的生成式智能体能否在虚拟社区中重现抑郁障碍的动态过程——以及心理治疗的对照试验能否完全在 \emph{in silico} 中运行。
先前的生成式智能体仿真运行在以天计的时间尺度上，且不带临床表型；而认知行为疗法（CBT）的真实随机对照试验（RCT）成本高昂、历时数月，且很少支持多于少数几个平行臂。
本文研究基于生成式智能体的抑郁症社区仿真，目标有二：(1) 虚拟社区中涌现的、有临床依据的症状动态；(2) 作为方法学工具的 \emph{in silico} 多臂 CBT 试验。
具体地，我们提出 \sys——在这一框架中，患者智能体由四层提示词栈、三步验证主诉图和动态抑郁状态机构建，生活在由居民智能体组成的村庄中，历时 36 个仿真日。
七臂对照协议（完整 CBT、无干预对照、三个效价化社会接触臂、支持性咨询和记忆写入消融）在冻结快照上、以 PHQ-9、BDI-II 和超短量表、在双层复制设计下评估。
大量实验表明：仿真的 CBT 使抑郁症状的减轻超出自然病程（\ph{X.X} 对 \ph{X.X} 个 PHQ-9 分），将认知重构与非特异性支持分离开来（较支持性咨询 +\ph{X.X} 分），并且阻断记忆写入保留急性反应，其持久性读数（\ph{XX}\% 对 \ph{XX}\% 反弹）已在随访扩展下登记——目前尚无同版本随访数据——且相对无干预对照在 PHQ-9 降幅上的增益最大（+\ph{XX}\%）。
所有数值均为占位符，待目标批次核验。
代码与配置可在 \ph{URL} 获取。
\end{abstract}

%%====================================================================
%%  1  引言
%%====================================================================
\section{引言}\label{sec:intro}

抑郁症是全球负担最重的精神障碍之一，估计影响 2.8 亿人~\citep{who2022}。
认知行为疗法（CBT）是一线、指南推荐的治疗，荟萃分析效应量为中到大~\citep{cuijpers2013,hofmann2012,cuijpers2023}。
然而结局仍然具有异质性，CBT 特定技术相对于联盟、接触与共情倾听等非特异性因素的独特贡献长期存在争议~\citep{lambert1992}，而维持治疗结束后收益的过程无法在真实患者身上进行实验性操纵。

生成式智能体学会在模拟社会中表现出令人信服的行为~\citep{park2023}。
这引出一个自然的问题：生成式智能体能否同样做到——在模拟社区中恶化、应对、求助并对心理治疗产生反应，从而使对照治疗试验可以完全在 \emph{in silico} 中运行？

先前的生成式智能体社会尚未达到这一情形。
原始的生成式智能体小镇运行在两天的时间尺度上，且仅模拟健康的社会行为~\citep{park2023}。
OASIS 将社交媒体互动扩展到一百万智能体，但仍停留在天与帖子的尺度~\citep{oasis2024}。
Agentopia 将生命仿真扩展到十个仿真年，展示了涌现的社会动态与学习~\citep{agentopia2026}，但其智能体过着健康的生活：没有任何社会嵌入精神表型、提供治疗协议或测量标准化临床结局。
将生成式智能体用于心理健康决定因素的提案仍是概念性的~\citep{kambeitz2025}。

临床一侧存在镜像的缺口。
真实 CBT 试验成本高昂、历时数月至数年，且很少支持多于少数几个平行臂；不提供循证治疗受伦理约束，因此成分分离设计难以实施。
记忆巩固——治疗性经验的编码与复述——在人类身上无法在不带自身混淆因素的范式之外被实验性阻断。
计算精神病学提供机制精度~\citep{montague2012,hauser2022,friston2023,zavlis2025}，但很少呈现治疗对话——心理治疗实际运作的媒介。

本文提出 \sys，一个生成式智能体平台，它将临床级别的抑郁表型嵌入 Stanford-Town 风格的虚拟社区，并在其中运行多臂 CBT 对照试验。
图~\ref{fig:overview} 给出总览。
(1) \sys 通过四层提示词栈、一张每个转移都必须通过三步验证管线的主诉图，以及一个动态抑郁状态机，把抑郁表型嵌入生成式智能体，避免了先前智能体社会的无症状仿真。
(2) \sys 运行受控的 \emph{in silico} 试验：在匹配接触剂量下、36 个仿真日内同步对照七个臂，这一臂数超出了真实 CBT 试验通常随机化的规模。
(3) \sys 使每一条声明都可审计：冻结快照、双层复制和盲态证据审计把每一个已提交的主诉图转移逐字锚定到患者原话上。

仿真结构为一个周期性、分阶段的循环。
一次运行经过 72 个仿真步推进，每步 720 仿真分钟，即 36 个仿真日。
调度器每 6 步插入一次受控接触——医生会谈（阶段 1）或居民聊天（阶段 2）。
对话与事件写回智能体的记忆流；动态抑郁状态机逐轮更新；主诉图仅在被验证的最小变化提交之后推进。
在七个评估节点，世界冻结快照，由独立的复评管线施测标准化量表。

我们以跨两个人设库的 \ph{N} 个人设、七个试验臂、轨迹与相关分析、案例研究以及骨干模型比较来评估 \sys，总结于表~\ref{tab:comparison}。
仿真的 CBT 使 PHQ-9 降低 \ph{X.X} 分，而无干预对照为 \ph{X.X} 分（相对改善 +\ph{XX}\%）；阻断记忆写入保留急性反应，其持久性读数（\ph{XX}\% 对 \ph{XX}\% 反弹）已在随访扩展下登记。
本文所有定量结果均为占位符，待目标批次核验；我们对其显式标注，且绝不将其解读为临床效应量。

我们的贡献总结如下：
\begin{itemize}
\item \textbf{系统。}据我们所知，\sys 是第一个嵌入临床级别抑郁表型（四层提示词栈、经验证主诉图、动态抑郁状态机）并在仿真回路内提供结构化 12 次会谈 CBT 协议的生成式智能体社区。
\item \textbf{工具。}我们将 \emph{in silico} 多臂试验操作化：匹配接触剂量下的七个臂——完整 CBT、无干预对照、三个效价化社会接触、支持性咨询和记忆写入消融——配以冻结快照心理测量（PHQ-9、BDI-II、超短量表）与双层复制。
\item \textbf{分离与审计。}该试验将特定技术与非特异性支持分离开来（G1 对 G6），并探究记忆巩固在持久性中的作用（G7）；同时，PSI-Bench 风格的拟人度检查、人设恢复和盲态证据审计把每个状态转移锚定在可审计的患者证据上。
\end{itemize}

%%====================================================================
%%  2  相关工作
%%====================================================================
\section{相关工作}\label{sec:related}

\subsection{基于智能体的社会仿真与生成式智能体}\label{sec:related-abm}

近期工作探索了将基于智能体的模型作为模拟人类系统的工具，从参数化的基于规则的智能体~\citep{bonabeau2002} 到 LLM 驱动的社会。
\citet{park2023} 在一个小镇中引入了具有记忆、反思与规划的生成式智能体，在两个仿真日内产生可信的社会行为。
OASIS 将这一范式扩展到社交媒体平台上的一百万智能体~\citep{oasis2024}。
Agentopia 将生命仿真扩展到 100 个智能体、十个仿真年，展示了涌现的社会动态与仿真内学习~\citep{agentopia2026}。
\citet{kambeitz2025} 提出将生成式智能体作为心理健康的环境与社会决定因素的建模工具，但该提案仍是概念性的。
这些社会在各种尺度上模拟健康的社会生活；没有一个嵌入精神表型，没有一个提供治疗协议，也没有一个测量标准化临床结局——有临床依据的社区仿真因此仍属空白。

\subsection{LLM 角色扮演、模拟来访者与计算精神病学}\label{sec:related-clinical}

近期工作探索了 LLM 角色扮演及其认识论边界~\citep{shanahan2023}、LLM 在精神病学中的应用~\citep{volkmer2024,stade2024}，以及有理论依据的心理健康来访者仿真：PatientAct 将模拟来访者锚定于临床理论并建模动态内部状态，并在咨询师训练对话中对其评估~\citep{patientact2026}。
在治疗一侧，荟萃分析确立了 CBT 的疗效及其异质、机制存争议的性质~\citep{cuijpers2013,hofmann2012,cuijpers2023,lambert1992}。
计算精神病学构建精神障碍的机制模型~\citep{montague2012,hauser2022,friston2023,zavlis2025}，但这些模型很少以自然语言呈现治疗对话。
模拟来访者生活在单一的临床对话中，真实试验昂贵且臂数少，二者都没有运行在一个症状、社会接触与治疗跨周交互的社会社区中——生活社区内的多臂 \emph{in silico} 试验因此仍属空白。

%%====================================================================
%%  3  MoodTown
%%====================================================================
\section{\sys}\label{sec:system}

\sys 在 Stanford-Town 生成式智能体框架~\citep{park2023} 之上扩展了抑郁表型、CBT 治疗管线和试验级评估层。
该系统是规则编排与 LLM 生成的混合体。
规则负责时间步进、调度、条件分配、会谈推进和快照触发。
LLM 负责日程规划、自然语言交互、反思、症状状态更新和主诉图转移。
图~\ref{fig:overview} 给出总览；完整设计细节见附录~\ref{app:details}。

%% ---------------------------------------------------------------
\subsection{智能体设计}\label{sec:agent}

\sys 中每个患者智能体由四个协同组件组装而成。
完整提示词见附录~\ref{app:prompts}。

\paragraph{档案。}
基础档案承载姓名、年龄、特质、生活习惯、每日计划和当日情境。
这些属性在一次运行中相对稳定，每个仿真日刷新一次。

\paragraph{临床病例锚点。}
每个患者承载一份长期病例背景、一个根主诉锚点和一个遵循 Beck 抑郁认知模型~\citep{beck1967,beck1979} 的固定核心信念。
这些属性在整个运行期间保持固定：主诉图管线不能改写它们，触及它们的 LLM 提案会被代码级门控拒绝。
附着于当前主诉阶段的近期相关经历在每次授权会谈后更新。

\paragraph{动态症状状态。}
两个耦合组件跟踪症状。
其一，动态抑郁状态机——改编自 PatientAct 有理论依据的动态来访者状态设计~\citep{patientact2026}——从每日事件和对话演化，在包括无干预对照在内的每个臂中以相同方式运行，并驱动下游行为。
其二，主诉图持有当前主诉阶段：一个至多 100 字符的标签和至多两句的摘要（一句关于持续存在的困难，一句关于最近的最小变化及其局限）。
状态机逐轮更新，而主诉阶段仅在被验证的最小变化提交之后改变——这在设计上是自然的，因为认知应当比情绪变化得更慢。

\paragraph{记忆与社会支持。}
每个智能体维护一个由事件、聊天和思想节点组成的本地记忆流，可按新近性、重要性和相关性检索。
初始化时每个患者接收四条种子压力记忆：一个压力事件、一个羞耻触发、一个消极自我念头和一个情境性紧张。
咨询历史模块检索并汇总既往会谈，使治疗跨周保持连续。
社会支持网络由治疗师智能体、三个脚本化居民接触智能体和旁观村民组成。
记忆写入消融臂（G7）阻断所有患者事件/思想/聊天写入，而其他一切照常运行。

所有四个组件通过四层提示词栈渲染进 LLM 上下文：(1) 基础人格，(2) 当前主诉阶段加稳定锚点，(3) 会话上下文（地点、时间、对象、关系、互动类型），(4) 瞬时情绪与表达引导（标签、风格、强度、0--10 量表上的袒露与防御、相对上一轮的波动、自我修正模式）。
第 (4) 层约束单轮表达，不得被解读为持续的症状变化。
完整的层内容见附录~\ref{app:details}。

%% ---------------------------------------------------------------
\subsection{仿真流程}\label{sec:procedure}

一次 \sys 运行分三个阶段进行。
初始化构建世界、分配人设及其锚点、播种四条压力记忆，并在记忆播种之后、患者第一次自主生活决策之前捕获 T0 基线快照。
主循环随后推进 72 个仿真步、每步 720 仿真分钟——合计 36 个仿真日。
每步引出当日计划、在村庄地图上执行动作与交互、插入规则触发的受控接触、把对话与事件写回记忆、更新动态抑郁状态，并在评估节点冻结快照。
因此，12 次会谈时代的一次运行安排 12 次受控接触和 7 个评估节点（T0、S2、S4、S6、S8、S10、S12）。
无计划接触的臂（无干预对照与自然接触臂）接收按步数对齐、带虚拟会谈标签的量表节点，这些标签表示时钟对齐而非实际交付的会谈。
完整的循环细节见附录~\ref{app:details}。

%% ---------------------------------------------------------------
\subsection{接触与调度}\label{sec:contact}

受控接触是干预到达患者的途径。
调度器每 6 个仿真步触发一次，预订两种阶段之一。
阶段 1 预订与治疗师智能体的一次 2{,}881 仿真分钟的医生会谈；阶段 2 预订与脚本化居民的一次 1{,}440 仿真分钟的居民聊天。
调度完全由规则驱动：接触被插入每日计划，绝不交给智能体的自发性，因此剂量由构造控制。
交付 G1 型受控会谈的臂运行强制聊天（forced-chat）记忆策略，使被调度的对话被记作治疗而非日常寒暄。
自然接触臂（G10、G11）从不调度；其效价化居民前缀仅在居民于自然对话中向患者搭话时注入。
咨询室场景（G4）将地图与角色重配为两个角色，同时继承同一基础 CBT 链与控制器。
完整的调度机制见附录~\ref{app:details}。

%% ---------------------------------------------------------------
\subsection{七臂 \emph{in silico} 试验协议}\label{sec:arms}

\sys 在单一村庄配置内操作化一个七臂对照试验：每个臂是深合并到同一基础配置之上的组配置叠加，因此各臂共享世界、人设库和调度器，仅在干预叠加上有差异。
附录~\ref{app:settings} 中的表~\ref{tab:registry} 列出包括扩展条件与非激活条件在内的完整组注册表。

各臂如下：
\begin{itemize}
\item \textbf{G1 —— 完整 CBT。}基准臂：12 次会谈结构化 CBT，带会谈提示词、对话判断、会谈评估、咨询历史、带环境反馈的作业提取、动态抑郁、记忆注入和强制聊天记忆策略。
\item \textbf{G2 —— 无干预对照（自然病程）。}无干预：会面、会谈提示词、判断和咨询历史全部关闭；抑郁状态机与记忆继续运行，量表节点带虚拟标签。
\item \textbf{G3 —— 中性社会接触。}十二次中性效价的计划居民聊天：仅日常生活闲谈——不安慰、不分析、不建议。
\item \textbf{G5 —— 负面社会接触。}与 G3 相同的日程，加负面效价前缀：低回应性、将痛苦最小化、指责、否认努力与需求。
\item \textbf{G9 —— 积极社会接触。}相同日程，加积极效价前缀：简短而熟悉的温暖、具体可信的肯定、愿意陪伴——没有空洞鼓励、诊断或强加建议。
\item \textbf{G6 —— 支持性咨询。}与 G1 相同的接触剂量，但关闭 CBT 特有机制（无会谈提示词、无判断、无会谈评估、无作业）：一个用于分离非特异性支持的主动治疗对照。
\item \textbf{G7 —— 记忆写入消融。}与 G1 完全相同的 CBT，但阻断所有患者事件/思想/聊天记忆写入；动态抑郁状态机仍然运行，因此这是记忆写入消融，而非状态消融。
\end{itemize}

CBT 疗程是一个固定的 12 次会谈大纲（会谈 1 $\rightarrow$ 2.1 $\rightarrow$ 2.2 $\rightarrow$ 2.3 $\rightarrow$ 3.1 $\rightarrow$ 3.2 $\rightarrow$ 3.3-A $\rightarrow$ 3.3-B $\rightarrow$ 4.1 $\rightarrow$ 4.2 $\rightarrow$ 4.3 $\rightarrow$ 4.4），历经签到、跨认知三角的自动化思维识别、认知歪曲标记、苏格拉底式重构、行为激活和复发预防~\citep{beck1979}。
在每次会谈内部，规划被拆分到本地模块：会谈任务规划器（Session Task Planner）维护具有单调进展的小目标（无 $\rightarrow$ 部分 $\rightarrow$ 完成）；选择器（Selector）仅从受控候选集中挑选回应策略与微技能；一个轻量判断器保留轻量判断加会谈终止，终止被门控——仅当小目标进展达到完成阈值或会谈达到估计 40 分钟之后，常规结束判断才开放。
信息边界始终成立：判断器与治疗师接收长期病例背景和可观察的结构化状态，绝不接收患者的内部变量。
三级风险路由器（无 / 可能 / 高）作为安全机制门控会谈流程。
我们的协议库中设计了六个治疗家族（CBT、IPT、BA、DBT、PST、SPT）；主线实现 CBT。
完整的会谈机制与阈值见附录~\ref{app:details}。

%% ---------------------------------------------------------------
\subsection{环境设计}\label{sec:environment}

\paragraph{环境模型。}
村庄提供地图、地点和居民生态，包括治疗师（蜻蜓队长）、三个脚本化接触居民和旁观村民。
日程规划、移动和自然对话都发生在这张地图上，因此治疗被嵌入持续进行的社会生活之中，而非一间空诊室。

\paragraph{作业反馈回路。}
每次 G1 型会谈后，系统提取至多三项作业，在环境模型中模拟其结果，并把结果作为记忆注入——闭合治疗与生活之间的回路。

\paragraph{试验编排。}
批量编排器将每个组配置深合并到基础配置上，运行仿真，冻结快照，并以来源信息（组、人设、严重度、控制器身份、内容哈希）归档对话、量表、记忆和主诉图证据。
不同控制器身份下的运行绝不混池；控制器身份是一个分层变量。

%%====================================================================
%%  4  结局与评估
%%====================================================================
\section{结局与评估}\label{sec:outcome}

\subsection{结局指标}\label{sec:outcomes}

\sys 以冻结快照上的阶段性心理测量来度量症状。
在两个终点节点（T0 和 S12），每次运行施测 PHQ-9（9 项，每项 0--3）~\citep{kroenke2001} 和 BDI-II（21 项）~\citep{beck1996}，每个快照各施测 $K=10$ 次独立施测。
在五个中间节点（S2--S10），每次运行施测两个 5 项超短量表——总体抑郁水平与功能损害，以及总体焦虑水平与功能损害（每项 0--4，总分 0--20）——各 $K=5$。
因此一次运行共安排 90 次量表施测（$2 \times 10 + 2 \times 10 + 5 \times 2 \times 5$）。

一个臂的主要结局是其各次运行的 PHQ-9 平均变化，
\begin{equation}\label{eq:delta}
\Delta_i \;=\; x_i(\mathrm{T0}) \;-\; x_i(\mathrm{S12}),
\end{equation}
其中更大的 $\Delta$ 表示更大的症状减轻，以及其相对形式
\begin{equation}\label{eq:deltapct}
\Delta\%_i \;=\; 100 \cdot \Delta_i / x_i(\mathrm{T0}).
\end{equation}
测量采用冻结快照协议。
仿真期间，世界仅冻结可恢复的运行状态（配置、触发元数据、智能体检查点、对话与抑郁状态、内容哈希、来源信息）；随后由独立的复评管线从归档中施测量表。
每一项都在单轮临时对话中作答：不追加聊天记忆，不提交动态状态，后续条目看不到先前答案——这是一个无状态协议，不同于连续的临床访谈，我们对此显式指出。
长量表由专家评分 LLM 根据自然语言答案逐项评分（0--3）；短量表采用规则优先评分，直接提取优先，而 LLM 兜底的输出会被任何可直接提取的分数覆盖，留下分数来源审计轨迹。
规则校验器复查条目数（9 / 21 / 5）、取值范围、总和一致性、严重度分类以及 PHQ-9 第 9 项安全字段一致性。
随访扩展（FU30--FU120）用于持久性读数；它在 120 步批次族中启用，而当前默认的 36 天窗口止于 S12，我们据此对持久性结果作相应标注。
完整的指标定义见附录~\ref{app:settings}。

\subsection{模拟测量的信度}\label{sec:reliability}

评估在设计上是双层的。
外层复制独立的仿真运行（$R01, R02, \dots$），回答轨迹能否重现。
内层在同一冻结快照上重复量表生成 $K$ 次，回答测量是否稳定。
两层绝不合并为一个样本量；比较的单位是独立运行。

我们报告快照内样本标准差与 95\% $t$ 区间、在单向随机效应组内相关系数下的量表水平重测信度~\citep{shrout1979}
\begin{equation}\label{eq:icc1}
\mathrm{ICC}(1,1) \;=\; \frac{\mathrm{MS}_R - \mathrm{MS}_E}{\mathrm{MS}_R + (K-1)\,\mathrm{MS}_E},
\qquad
\mathrm{ICC}(1,K) \;=\; \frac{\mathrm{MS}_R - \mathrm{MS}_E}{\mathrm{MS}_R},
\end{equation}
其中 $\mathrm{MS}_R$ 与 $\mathrm{MS}_E$ 为目标间与目标内均方；条目水平一致性用二次加权 kappa~\citep{cohen1968}
\begin{equation}\label{eq:kappa}
\kappa_w \;=\; 1 - \frac{\sum_{i,j} w_{ij}\, p^{o}_{ij}}{\sum_{i,j} w_{ij}\, p^{e}_{ij}},
\qquad
w_{ij} \;=\; \left( \frac{i-j}{m-1} \right)^{2},
\end{equation}
并以同节点总分相关和前后测变化相关来度量 PHQ-9/BDI-II 趋同。
效应量为 Cohen's $d$，附 95\% 置信区间~\citep{cohen1988}；组间对比使用以人设和运行为随机效应的线性混合模型，并对多重比较作 Bonferroni 校正。

\subsection{拟人度与证据审计}\label{sec:humaneval}

三个审计族把仿真锚定于证据。
其一，PSI-Bench 风格的话语标注为每条患者回应标注一个 PTC 类别（痛苦 / 探索--反思 / 领悟--接纳 / 填充）和九种情绪之一；标注仅使用每次会谈的前 8 条患者回应，PTC 至多看到其前 6 条医患消息、情绪至多 4 条，绝看不到未来轮次。
其二，静态人设恢复检验智能体的话语能否在四选一强迫选择中（对三个干扰项）恢复其三个锚点（根主诉、核心信念、长期背景）。
其三，证据审计复查每一个已提交的主诉图转移：一个对图更新、反思、量表和臂分配全盲的离线判断器，仅看到逐字患者话语和活动归档，并判断已提交的变化是否有证据支持。
我们将 PSI-Bench 读数报告为四个语料库、每个 64 条（8 次会谈 $\times$ 8 条回应；共 256 条）；第五个已下载语料库在两张已报告表格中均缺席，等待归档确认；逐轮 JSD 距离尚未计算；我们也未获得真实人类参照语料库——因此我们不作人类平价声明。

%%====================================================================
%%  5  实验
%%====================================================================
\section{实验}\label{sec:exp}

\subsection{实验设置}\label{sec:setup}

完整细节见附录~\ref{app:settings}。

\paragraph{社区数据。}
\begin{itemize}
\item \textbf{村庄（默认）。}完整小镇社区，含患者、治疗师、三个脚本化接触居民和 \ph{N} 名旁观村民。它凸显症状、社会接触与治疗如何在生态丰富的环境中交互。
\item \textbf{咨询室（G4）。}双角色场景，剥离社会环境而保留 CBT 链与控制器。它凸显治疗效应有多少能在没有社区的情况下存活。
\item \textbf{选择器人设库。}八个人设，年龄 18--66，覆盖校园、职场、家庭与晚年应激源（表~\ref{tab:personas}）。它凸显在匹配严重度下轨迹如何跨生活领域变化。
\item \textbf{KBD 人设库。}九个人设，各配四种严重度配置——一个未分级基础配置加轻度 / 中度 / 重度，共 36 个配置；运行时严重度开关仅在三个分级档位之间切换。它凸显跨严重度的剂量--反应结构，并锚定当前的中度默认设置。
\end{itemize}

\paragraph{仿真。}
每次运行将一名患者智能体嵌入村庄，共 72 步、每步 720 分钟（36 个仿真日）、12 次受控接触和 7 个评估节点，每次运行产生 90 次量表施测。
七臂试验覆盖 $7 \times \ph{R}$ 次运行——复制次数按批次设定（近期批次：$\times 1$ 或 $\times 2$；整个项目的范围：1--5）——共产生 \ph{N,NNN} 次阶段性评估和 \ph{N,NNN} 次量表施测，另有 \ph{N,NNN} 个独立复评的冻结快照。
控制器身份或提示词版本不同的批次绝不混池。

\paragraph{模型。}
医生、患者与评估链运行于 DeepSeek。
本地 Qwen3-8B 充当思考模型（会谈任务规划器、选择器、汇总与压缩），并充当证据审计的离线判断器。
检索使用 BGE-M3 嵌入；计算集群在 4 张 GPU 上部署 3 个 Qwen 服务和 3 个 BGE 服务。
思考模式仅限医患对话——在参考批次中，判断器调用的 token 约有 97.5\% 为推理 token。
模型分配属于运行配置，按批次清单报告，并非方法的固定属性；PSI-Bench 标注在本地思考模型（Qwen3-8B）上离线运行。

\paragraph{指标。}
疗效指标（$\Delta$、$\Delta\%$、轨迹、条目级变化、人设分层）、信度指标（样本 SD、95\% $t$ 置信区间、ICC$(1,1)$ / ICC$(1,K)$、二次加权 $\kappa_w$、PHQ-9/BDI-II 趋同），以及拟人度指标（PTC、情绪、人设恢复、证据审计通过率）。
完整的指标定义见附录~\ref{app:settings}。

%% ---------------------------------------------------------------
\subsection{长期社区仿真分析}\label{sec:simanalysis}

图~\ref{fig:traj} 绘出 36 天窗口内的症状轨迹，图~\ref{fig:heatmap} 将行为频率与症状变化相关联。
我们突出四个发现：
(1) \textbf{自然病程几乎平坦。}无干预轨迹在 36 天内变化 \ph{X.X} 个 PHQ-9 分，有波动但无结构性恢复；这表明抑郁状态机在无干预下维持症状，而这正是自然病程对照所要求的。
(2) \textbf{不良社会接触主动恶化症状。}负面接触（G5）相对无干预对照使 PHQ-9 升高 \ph{X.X} 分，映照现实世界的人际应激与社会节律紊乱——社会逆境是加剧症状，而不仅仅是未能改善症状。
(3) \textbf{行为与症状耦合。}作业完成度与 PHQ-9 改善的相关为 $r = \ph{0.XX}$，社交活动频率的相关为 $r = \ph{0.XX}$（图~\ref{fig:heatmap}）；这提示一种与真实 CBT 中报告的作业剂量效应~\citep{burns2001} 相似的投入剂量结构。
(4) \textbf{认知变化是分阶段的，而非突变的。}已提交的主诉图转移集中于重构阶段（第 \ph{X}--\ph{Y} 次会谈中的转移占 \ph{XX}\%），且 \ph{XX}\% 的已提交转移通过盲态证据审计；这表明症状变化依托于经验证的最小认知变化——鉴于三步验证管线，这在设计上是自然的。

%% ---------------------------------------------------------------
\subsection{七臂 \emph{in silico} 试验}\label{sec:trial}

表~\ref{tab:main} 报告主试验读数：基线（T0）、疗程中段（S6）与终点（S12）的 PHQ-9，附变化列与相对无干预对照的效应量，并附一行真实世界荟萃分析锚点~\citep{cuijpers2013}。
图~\ref{fig:traj} 展示完整轨迹。

\paragraph{\emph{In silico} CBT 使抑郁症状的减轻超出自发缓解。}
G1 使 PHQ-9 改善 \ph{X.X} 分（较基线 \ph{XX}\%，相对 G2 的 $d = \ph{X.XX}$），最陡的改善出现在第 \ph{X} 次与第 \ph{Y} 次会谈之间；12 次会谈大纲的认知单元位于第 2--4 次会谈，因此把最陡区段归因于某一特定大纲阶段本身就是一项读数。
PHQ-9 与 BDI-II 趋同于 $r = \ph{0.XX}$，支持跨工具一致性。

\paragraph{在匹配接触下，结构化 CBT 优于支持性咨询。}
G1 超出 G6 \ph{X.X} 个 PHQ-9 分（$d = \ph{X.XX}$，95\% CI \ph{X.X, X.X}，Bonferroni 校正后 $p = \ph{X.XX}$），该优势在各人设间一致（Group$\times$Persona $p = \ph{X.X}$）。
两臂共享完全相同的接触日程，因此这一差异在治疗联盟、定期接触与共情倾听之上分离出了认知重构——歪曲识别、证据检验、行为实验。

\paragraph{社会接触效价产生分级的症状梯度。}
终点排序为 G9 $>$ G3 $>$ G2 $>$ G5：积极接触改善（$\Delta = $ \ph{X.X}），中性接触与无干预对照无法区分（$d = \ph{X.XX}$，$p = \ph{X.X}$），负面接触恶化症状（$\Delta = -$\ph{X.X}）。
中性接触的零结果是一项结果，而非失败：匹配剂量下的单纯陪伴在自然病程之外没有额外增益，提示缺乏治疗结构的非特异性社会接触是惰性的，而效应由效价承载。

\paragraph{阻断记忆写入保留急性反应；持久性是已登记的读数。}
G7 在 S12 与 G1 相当（$d = \ph{X.XX}$，$p = \ph{X.X}$）：在没有新记忆形成的情况下急性反应得以保留，而这正是当前 36 天窗口所支持的读数。
持久性预测是登记在案的而非已观察的：随访节点仅存在于更早的 120 步批次族（\S\ref{sec:outcomes}），且尚无同版本的 G7 随访数据，因此我们把反弹重入（\ph{XX}\% 对 \ph{XX}\%，$\chi^2 = \ph{X.X}$，Group$\times$Time $F = \ph{X.X}$）与更少的向残余/恢复阶段的转移（图~\ref{fig:graph}）登记为随访扩展下的计划分析。
若获证实，这一分离——振幅保留、持久性改变——将与记忆巩固在维持收益中的仿真内因果角色相一致；它是否延伸至人类记忆巩固，则是一个有待临床转化的假设，而非声明。

\begin{table}[t]
\centering
\caption{七臂 \emph{in silico} 试验：各臂 PHQ-9。
$\Delta = \mathrm{T0} - \mathrm{S12}$（式~\ref{eq:delta}）；$\downarrow$ 表示症状减轻（改善），$\uparrow$ 表示恶化；$d$ 为相对 G2 的 Cohen's $d$；锚点行引用真实世界 CBT 对等待列表的荟萃分析效应~\citep{cuijpers2013} 作为量级参照，而非等效声明。
加粗 = 每列最高、下划线 = 每列最低，将在数据锁定时套用。
所有数值均为占位符，待目标批次核验。}
\label{tab:main}
\small
\begin{tabular}{@{}llcccccc@{}}
\toprule
臂别 & 代码 & T0 & S6 & S12 & $\Delta$ & $\Delta\%$ & 对 G2 的 $d$ \\
\midrule
完整 CBT               & G1 & \ph{X.X} & \ph{X.X} & \ph{X.X} & \ph{X.X}~$\downarrow$ & \ph{XX}\% & \ph{X.XX} \\
支持性咨询 & G6 & \ph{X.X} & \ph{X.X} & \ph{X.X} & \ph{X.X}~$\downarrow$ & \ph{XX}\% & \ph{X.XX} \\
记忆写入消融  & G7 & \ph{X.X} & \ph{X.X} & \ph{X.X} & \ph{X.X}~$\downarrow$ & \ph{XX}\% & \ph{X.XX} \\
积极接触       & G9 & \ph{X.X} & \ph{X.X} & \ph{X.X} & \ph{X.X}~$\downarrow$ & \ph{XX}\% & \ph{X.XX} \\
中性接触        & G3 & \ph{X.X} & \ph{X.X} & \ph{X.X} & \ph{X.X}~$\downarrow$ & \ph{XX}\% & \ph{X.XX} \\
无干预对照  & G2 & \ph{X.X} & \ph{X.X} & \ph{X.X} & \ph{X.X}~$\downarrow$ & \ph{XX}\% & --- \\
负面接触       & G5 & \ph{X.X} & \ph{X.X} & \ph{X.X} & $-$\ph{X.X}~$\uparrow$ & \ph{XX}\% & \ph{X.XX} \\
\midrule
\multicolumn{8}{@{}l}{\emph{锚点}：真实世界 CBT 对等待列表，荟萃分析 $g = \ph{X.XX}$~\citep{cuijpers2013}（参照行）。} \\
\bottomrule
\end{tabular}
\end{table}

我们按主题解读表~\ref{tab:main}。
\begin{enumerate}
\item \textbf{总体疗效。}G1 改善 \ph{X.X} 分，而无干预对照为 \ph{X.X} 分；这表明完整链条——会谈结构、判断、作业回路、记忆——而非时间本身驱动改善。
\item \textbf{特定技术。}在匹配接触下 G1 超出 G6 \ph{X.X} 分；这在非特异性支持之上分离出认知重构这一有效成分，呼应围绕功能失调态度改变的机制争论~\citep{burns2001,lambert1992}。
\item \textbf{社会效价。}G9 $>$ G3 $\approx$ G2 $>$ G5 的排序提示社会环境方向性地调节症状，映照真实抑郁症中的人际应激发现。
\item \textbf{记忆巩固。}G7 保留的终点提示急性振幅并不需要新记忆形成；其持久性读数（\ph{XX}\% 反弹）已在随访扩展下登记，但尚未在同版本批次中观察到，因此振幅--持久性分离仍是一个预测——会谈绑定的加工对已巩固的记忆。
\item \textbf{零结果与权衡。}G3 的零结果与 G7 保留的急性反应被作为发现报告：前者界定了单纯陪伴所能达成的范围，后者展示了记忆写入\emph{并非}为哪些目的所必需——二者都是有信息量的权衡，而非噪声。
\end{enumerate}

除七臂之外，设计还支持真实试验很少能实施的对比：场景消融（G4）移除社区，自然发生变体（G10、G11）用自发接触取代调度，模块消融移除单个架构组件。
场景对比与自然发生对比已经运行；模块消融已计划且尚未执行；附录~\ref{app:additional} 报告前者并登记后者。
单一冻结社区配置内的七个同步对照臂，其臂数超出了真实 CBT 试验通常随机化的规模；正是这一规模使该试验在方法学上有价值（表~\ref{tab:comparison}）。

%% ---------------------------------------------------------------
\subsection{计算成本}\label{sec:cost}

表~\ref{tab:cost} 报告每次运行的成本，图~\ref{fig:cost} 给出其阶段分解；数值为占位符，待目标批次。
成本由医患对话主导：思考模式仅在该区段运行，且参考批次中判断器调用的 token 约有 97.5\% 为推理 token。
咨询室场景运行降成本标志（lite 非咨询处理）。
因此主导瓶颈是会谈对话，而非世界仿真或量表施测——这提示驱动成本的是臂扫描，而非观测密度。

\begin{table}[t]
\centering
\caption{每次运行的计算成本（占位符，待目标批次核验）。
M = 百万 token；K = 千次调用；h = 墙钟小时。}
\label{tab:cost}
\small
\begin{tabular}{@{}lcccc@{}}
\toprule
臂别 & Token 数（M） & LLM 调用数（K） & 冻结快照 & 墙钟时间（h） \\
\midrule
G1 & \ph{X.X} & \ph{X.X} & 7 & \ph{X.X} \\
G2 & \ph{X.X} & \ph{X.X} & 7 & \ph{X.X} \\
G3 / G5 / G9 & \ph{X.X} & \ph{X.X} & 7 & \ph{X.X} \\
G6 & \ph{X.X} & \ph{X.X} & 7 & \ph{X.X} \\
G7 & \ph{X.X} & \ph{X.X} & 7 & \ph{X.X} \\
\bottomrule
\end{tabular}
\end{table}

%%====================================================================
%%  6  结论
%%====================================================================
\section{结论}\label{sec:conclusion}

我们提出了 \sys，一个把临床级别抑郁表型嵌入虚拟社区并在其中运行多臂 CBT 试验的生成式智能体平台。
四层提示词栈、三步经验证主诉图、动态抑郁状态机，以及冻结快照上的阶段性心理测量，把一个活的社会环境与试验级测量在 36 个仿真日上耦合起来。

社区展现出持续、分阶段的症状动态：无干预轨迹保持平坦，不良社会接触主动恶化症状，认知变化经由经验证的最小步骤而非突变跳变推进。

七臂试验分离了真实设计难以分离的东西：结构化 CBT 使 PHQ-9 改善 \ph{X.X} 分，而无干预对照为 \ph{X.X} 分（相对改善 +\ph{XX}\%），超出匹配剂量支持性咨询 \ph{X.X} 分；阻断记忆写入保留急性反应，其持久性读数（\ph{XX}\% 对 \ph{XX}\% 反弹）已在随访扩展下登记。
\sys 是面向治疗设计的假设生成工具——不是临床效应量的来源——其完整评估套件（信度、拟人度、盲态证据审计）旨在使其最终报告的每一个数字都可审计。

%%====================================================================
%%  局限性
%%====================================================================
\section*{局限性}\label{sec:limitations}

\paragraph{表型效度。}
LLM 渲染的患者是拟像，不是患者：语言上的可信并不意味着现象学上的忠实。
我们以有理论依据的状态设计~\citep{patientact2026}、PSI-Bench 风格话语审计、人设恢复和盲态证据审计来缓解；真实人类参照语料库尚未获得，因此不作人类平价声明——这是一个开放挑战。

\paragraph{诊断效度。}
所有量表均为 LLM 模拟的自我报告，由专家评分 LLM 加规则校验评分，而非临床施测工具。
我们以双长量表、趋同检查、ICC 报告以及 \S\ref{sec:outcomes} 中声明的无状态逐项协议来缓解；与临床采集数据的校准仍待完成。

\paragraph{协议保真度。}
12 次会谈大纲是手册化 CBT~\citep{beck1979} 的结构化近似，由单一 LLM 治疗师交付，其微技能被限制在受控候选集内。
我们以固定大纲顺序、完成阈值和逐批次控制器身份来缓解；专家评定的保真度评分仍待完成。

\paragraph{数值校准。}
本文的每一个定量结果都是占位符，待目标批次核验；荟萃分析锚点行是量级参照，不是等效。
复制次数随批次变化（1--5），且跨控制器批次不可混池。

\paragraph{计算约束。}
每个臂需要完整的 36 天社区仿真与密集推理（思考模式仅限医患对话，在其中约占判断器 token 的 $\approx$97.5\%）。
$K=10$/$K=5$ 的内层重复限定了测量成本；在完全复制下扫掠严重度档位目前力有不逮。

%%====================================================================
%%  影响声明
%%====================================================================
\section*{影响声明}

\sys 是在受控 \emph{in silico} 条件下进行假设生成与治疗设计探索的研究工具。
它不是诊断设备，不是风险评估工具，也不是临床效应量的来源；模拟的痛苦是合成的，仿真内的安全相关内容由专用风险路由器和独立安全机制处理。
该平台不取代真实治疗或真实试验：\emph{in silico} 因果声明（例如记忆写入分离）指的是仿真自身的机制，必须经由临床研究加以转化。
我们承诺报告阴性与零结果、可审计的证据轨迹和占位符纪律——数字只有在批次核验之后才进入本稿——以便在该技术到达临床决策支持场景时，这一工具的输出仍可被审查。

%%====================================================================
%%  参考文献
%%====================================================================
\bibliographystyle{plainnat}
\bibliography{agentopia-references}

%%====================================================================
%%  附录
%%====================================================================
\appendix

%% ---------------------------------------------------------------
\section{社区与人设创建}\label{app:personas}

\subsection{管线概览}

人设创建分五步进行。
\textbf{第 1 步：设定选择。}选择社区模板（村庄或咨询室）和人设库。
\textbf{第 2 步：人设起草。}依据库规范起草每个人设的年龄、生活领域和促发应激源。
\textbf{第 3 步：锚点分配。}冻结长期背景、根主诉锚点、核心信念和一组可恢复活动（人设能够现实地重新投入的活动）。
\textbf{第 4 步：压力记忆播种。}播种四条压力记忆（压力事件、羞耻触发、消极自我念头、情境性紧张）。
\textbf{第 5 步：注册表登记。}将人设连同其严重度配置、锚点、可恢复活动和播种记忆登记进配置注册表（schema 2.1.0），使批次清单能把每次运行追溯回其人设配置。
组配置中的患者槽位默认为占位人设（Kabuda），由人设选择器在运行时替换；组机制不受该替换影响。

\subsection{人设库}

表~\ref{tab:personas} 列出八个选择器人设；KBD 库提供九个人设、各四种严重度配置——一个未分级基础配置加轻度 / 中度 / 重度——共 36 个配置（schema 2.1.0）。
运行时严重度开关在三个分级档位之间切换；未设置严重度开关时适用未分级基础配置。
当前的项目默认为中度严重度。

\begin{table}[h]
\centering
\caption{选择器人设库（中度配置；年龄跨度 18--66）。}
\label{tab:personas}
\small
\begin{tabular}{@{}llcl@{}}
\toprule
代码 & 人设 & 年龄 & 应激源领域 \\
\midrule
LRN & Lin Ruoning & 18 & 校园合作误会 \\
GC  & Gu Chen & 18 & 新校园中的陌生感 \\
CY  & Chen Yu & 25 & 一段关系结束；难以开始学习 \\
TW  & Tang Wan & 28 & 职场比较 \\
ZYH & Zhou Yuanhang & 34 & 家庭住房与开支压力 \\
SQL & Su Qinglan & 32 & 围绕独立选择的家庭冲突 \\
ZMY & Zhao Mingyuan & 60 & 退休后失去节奏 \\
XFH & Xu Fanghua & 66 & 独居；渴望陪伴 \\
\bottomrule
\end{tabular}
\end{table}

KBD 库涵盖失业与失败（KBD1）、长期比较与低自我价值（KBD2）、社交羞耻与回避（KBD3）、停工（KBD4）、照护耗竭（KBD5）、创作被拒（KBD6）、网购差评（KBD7），以及两个人设——其完整定义存于代码仓库（KBD8--KBD9）。

%% ---------------------------------------------------------------
\section{系统设计细节}\label{app:details}

\subsection{四层提示词栈}

第 1 层（基础人格）：姓名、年龄、特质、生活习惯、每日计划、当日情境。
第 2 层（主诉阶段）：阶段标签与摘要、近期相关经历，以及稳定背景——根主诉锚点、固定核心信念、表达风格与提示。
第 3 层（会话上下文）：地点、时间、对象与关系、互动类型、场景特有的附加判断。
第 4 层（瞬时情绪）：标签、风格、强度、袒露与防御（0--10）、相对上一轮的波动、自我修正模式；该层的输出是单轮约束，不能支撑持续症状改善的声明。
渲染还附加咨询历史摘要（跨会话累积）、智能体记忆列表（事件/聊天/思想；G7 中写入被阻断），以及一个滚动对话窗口——由进展时间估计、更早会话的摘要和最近五轮医患交流的逐字记录组成。

\subsection{动态抑郁状态机}

状态机改编自 PatientAct 有理论依据的动态来访者状态~\citep{patientact2026}，独立于主诉图，在每个臂中启用（显式启用或经基础配置继承），包括无干预对照。
它从每日事件和对话更新，并驱动下游行为。
更早的设计曾把七维持续症状状态直接注入提示词；经过一次内部消融，这一提示词注入区从渲染层移除，而状态机本身继续运行并驱动动态。

\subsection{各臂记忆策略}

G1 型受控会话运行强制聊天记忆策略，使被调度的对话被编码为治疗。
G3/G5/G9 对居民聊天禁用该策略。
每次 G1 会话后，系统提取至多三条作业指令，环境模型模拟其结果，结果作为记忆注入。
G6 与 G7 保持一般记忆注入启用（例如初始化压力记忆），但禁用 G1 运行的会话后人设特定（KBD）记忆注入；G7 另通过记忆写入控制阻断所有患者事件/思想/聊天写入（动态抑郁状态机仍运行）。
咨询历史检索汇总既往会谈以保证连续性。

\subsection{会话内部机制}

当前机制把原先集中式的判断器拆分为本地模块。
会谈任务规划器维护每次会谈的小目标，并做单调进展更新（无 $\rightarrow$ 部分 $\rightarrow$ 完成）。
选择器从路由器控制的候选集中选择回应策略与微技能，无需额外的重模型调用。
判断器仅保留轻量判断加终止，终止被门控：仅当小目标进展达到完成阈值（60\%）或会谈达到估计 40 分钟之后，常规结束判断才开放；在此之前，终止输出被归一化为假。
会谈目标 60 分钟，硬上限 60 轮，第 50 轮后生成巩固摘要，自第 55 轮起建议收尾。
会谈时长按每分钟对话 180 字符估计。
会话后本地完成审计包含用于会谈推进的三次会话 40\% 回退。
信息边界始终成立：判断器与治疗师看到长期病例背景和可观察的结构化状态，绝看不到患者内部变量。
控制器身份（Legacy、Minimal、Progressive D）按批次清单记录；跨控制器的同名组绝不混池。

\subsection{主诉图验证管线}

每个阶段转移经过三次带代码门控的独立 LLM 调用。
\textbf{(1) 变化检测器（Change Detector）。}仅给定该次会话获授权的逐字患者文本，判断相对当前阶段是否发生了足以影响后续表达或行为的最小变化；严格 JSON 输出（has\_new\_change / change / reason）；适用密集反例约束——计划不是执行，瞬时症状不是持续的能力下降，“知道……但仍然……”不得截断为前半句，限定词（“可能”、“仅此一次”、“但仍然”）必须保留；反思只有通过领悟才能证明变化，不能通过计划或想起的念头证明。
\textbf{(2) 阶段规划器（Stage Planner）。}把下一阶段构造成当前阶段加经验证的变化；标签 $\le$ 100 字符；摘要 $\le$ 两句（持续存在的困难；最小变化及其局限）；必须保留父主诉轴；不得输出叙事焦点、声明、标识符、版本和类型；不得展开疗效、机制和未来展望。
\textbf{(3) 阶段校验器（Stage Validator）。}一次独立调用返回变化保真度与阶段保真度（通过 / 失败 / 未知）；两者必须都通过，逐句检查每条新断言，阻断同义重复节点以及夸大的改善或恶化。
只有当三步全部成功且两个保真度都通过时，系统才提交——分配新节点标识符、添加边并更新位置；终末阶段不允许后继。
初始化时冻结：病例背景、根主诉锚点、核心信念。
由代码继承、不可被 LLM 提案改写：风格、情绪基线、关系参数。
不可变：逐字记录、观察、审计原始文本。
医生话语、反思结论和 CBT 进展绝不计为患者事实证据；只有正式的患者表达才能作为新阶段的依据，证据不足则不推进，且所有三步轨迹都归档。
更早的证据优先匹配机制已被这一分阶段管线取代；各批次按其清单中记录的机制解读。

\subsection{环境与编排细节}

村庄作为批量的默认完整小镇模式运行；咨询室模式将基础配置覆盖为双角色地图与角色阵容，并带降成本标志，同时继承基础 CBT 链与控制器，且作为场景消融而非内容臂来分析。
风险路由器按三级（无 / 可能 / 高）为每一轮评级并门控会谈流程。
冻结快照存储可恢复的运行状态：运行配置、触发元数据、智能体检查点与本地存储副本、对话与动态抑郁状态（含病例级核心信念）、内容哈希，以及组/人设/严重度/控制器来源信息。
复评从原始归档中自动选取实际的阶段性节点；T0 在压力记忆播种之后、患者第一次自主生活决策之前捕获。

\subsection{配置参数}

表~\ref{tab:config} 列出当前默认配置的主要参数。

\begin{table}[h]
\centering
\caption{主要配置参数（当前默认，12 次会谈时代）。}
\label{tab:config}
\small
\begin{tabular}{@{}lll@{}}
\toprule
参数 & 取值 & 作用域 \\
\midrule
仿真步长 & 720 分钟 & 世界时钟 \\
观测窗口 & 72 步 $=$ 36 天 & 每次运行 \\
受控接触节奏 & 每 6 步（每次运行 12 次） & 调度器 \\
医生会谈预订 & 2{,}881 仿真分钟 & 阶段 1 \\
居民聊天预订 & 1{,}440 仿真分钟 & 阶段 2 \\
CBT 疗程 & 12 次会谈，固定大纲 & G1/G7 \\
评估节点 & T0, S2, S4, S6, S8, S10, S12 & 所有臂 \\
长量表 / $K$ & PHQ-9、BDI-II；$K=10$ & T0, S12 \\
短量表 / $K$ & 2 $\times$ 5 项（0--20）；$K=5$ & S2--S10 \\
完成阈值 & 60\% & 会谈推进 \\
最早结束检查 & 估计 40 分钟 & 判断器门控 \\
建议时长 / 硬上限 & 60 分钟 / 60 轮 & 会话长度 \\
滚动窗口 & 最近 5 轮交流 $+$ 更早会话摘要 & 上下文 \\
时长估计 & 180 字符 / 分钟 & 上下文 \\
作业指令 & 每次会谈 $\le$ 3 条 & G1/G7 \\
压力记忆 & 每位患者 4 条 & 初始化 \\
严重度默认 & 中度 & 当前批次 \\
\bottomrule
\end{tabular}
\end{table}

%% ---------------------------------------------------------------
\section{实验设置与指标定义}\label{app:settings}

\subsection{批次注册表与效度标注}

表~\ref{tab:batches} 列出近期参考批次（内部注册表标识符）。
批次标识符是内部运行注册表 ID，为可追溯性而保留。
历史批次在被引用时带效度标注：$\dagger$ = 重构前时代主诉图推进在早期索引后停滞；$\ddagger$ = 早期社会接触批次中居民侧提示词误注入；$\S$ = 最早医生臂中的会话后记忆注入；$\P$ = 某批次的叙事文本滞后于其配置（以配置表为准）。
0922 试点按村庄模式对待，因为其运行时归档包含第三个智能体，表明是完整小镇而非双角色场景；这一解读等待人工归档确认。

\begin{table}[h]
\centering
\caption{近期参考批次（内部 ID）。复制按配置计；条件与人设按各清单注册。}
\label{tab:batches}
\small
\begin{tabular}{@{}lllll@{}}
\toprule
批次 & 条件 & 人设 & 复制 & 窗口 \\
\midrule
0908 & G1 / G2 / G5 & LRN; SQL (G1) & $\times 1$ & 72 步 \\
0915 & G1 / G2 / G5 & LRN; SQL, TW, ZYH (G1) & $\times 1$ & 72 步 \\
0922 & G1 & KBD1（中度） & $\times 2$ & 72 步 \\
\bottomrule
\end{tabular}
\end{table}

\subsection{组注册表}

表~\ref{tab:registry} 列出完整注册表，包括七臂试验之外的条件。
G8 是评估侧敏感性条件，在重置抑郁状态后复评已归档的 G1 快照；它绝不进入优越性排名。
G12（自然积极接触）已注册，有一次记录在案的运行，其结果从未被记录；我们将其排除出所有分析并显式注明这一缺口——因此自然发生族仅覆盖中性（G10）与负面（G11）接触，而积极接触仅存在于计划形式（G9）。

\begin{table}[h]
\centering
\caption{组注册表。七臂试验由 G1、G2、G3、G5、G6、G7、G9 组成；扩展对比与注册状态一并列出以求完整。}
\label{tab:registry}
\small
\begin{tabular}{@{}llp{8.6cm}@{}}
\toprule
代码 & 家族 & 定义 / 状态 \\
\midrule
G1 & 治疗 & 完整 CBT 基准（主实验臂）。 \\
G2 & 对照 & 无干预；虚拟标签量表节点。 \\
G3 & 社会接触 & 计划的中性居民聊天。 \\
G4 & 场景消融 & 咨询室双角色场景；非内容臂。 \\
G5 & 社会接触 & 计划的负面居民聊天。 \\
G6 & 治疗对照 & 匹配剂量下的支持性咨询。 \\
G7 & 机制消融 & CBT $+$ 患者记忆写入阻断。 \\
G8 & 评估侧 & 对已归档 G1 的敏感性复评；排除于排名之外。 \\
G9 & 社会接触 & 计划的积极居民聊天。 \\
G10 & 自然接触 & 自然发生的居民话语上的中性效价。 \\
G11 & 自然接触 & 自然发生的居民话语上的负面效价。 \\
G12 & 自然接触 & 积极效价，已注册；一次未记录的运行；已排除。 \\
\bottomrule
\end{tabular}
\end{table}

\subsection{指标定义}

\paragraph{疗效。}
按式~\ref{eq:delta}--\ref{eq:deltapct} 的 $\Delta$ 与 $\Delta\%$，在 PHQ-9（主要）、BDI-II（趋同）和两个超短量表（密集中期读数）上评估；条目级变化与人设分层为次要。
改善方向：分数下降。

\paragraph{信度。}
快照内样本 SD 与 95\% $t$ 置信区间；按式~\ref{eq:icc1} 的 ICC$(1,1)$ 与 ICC$(1,K)$；按式~\ref{eq:kappa} 的条目级二次加权 $\kappa_w$；PHQ-9/BDI-II 趋同以同节点总分相关和前后测变化相关度量。

\paragraph{拟人度。}
每次会谈前 8 条患者回应上的 PTC 类别与九情绪分布（窗口为 6 / 4 条最近消息；无未来上下文）；三个锚点上的静态人设恢复准确率（四选一强迫选择）；在对图更新、反思、量表和臂全盲的判断器下，已提交主诉图转移的证据审计通过率。

\paragraph{推断。}
线性混合模型（$\Delta \sim$ Group $+$ (1$|$Persona) $+$ (1$|$Run)）；纵向对比用 Group$\times$Time；Cohen's $d$ 附 95\% CI~\citep{cohen1988}；Bonferroni 校正。

%% ---------------------------------------------------------------
\section{附加结果}\label{app:additional}

本附录中所有数值均为占位符，待目标批次核验。

\paragraph{移除社区后治疗反应保持完整。}
咨询室场景（G4）在 S12 与村庄 G1 相当（$d = \ph{X.XX}$，$p = \ph{X.X}$）；这表明急性治疗效应不依赖社会环境，在急性窗口内界定了社区的贡献——其在持久性中的作用仍有待检验。

\paragraph{自然与计划接触存在差异。}
自然发生的中性接触（G10）使症状移动 \ph{X.X}，而计划的中性接触 G3 为 \ph{X.X}；G11 相对 G5 相差 \ph{X.X}；这提示时机与框架调节社会接触效应，映照偶然支持与安排支持之间的差异。

\paragraph{模块消融已登记，尚未执行。}
组注册表中没有移除主诉图、逐轮情绪推断或关系修饰器的臂，且迄今未运行任何模块消融批次。
已登记的计划包含约九次额外运行，每个对比一个模块。
计划读数：移除主诉图应使治疗反应趋平（$\Delta$PHQ-9 读数，计划 \ph{X.X}）；移除逐轮情绪推断应降低真实性检查（连续性 \ph{X.X} $\rightarrow$ 计划 \ph{X.X}）并削弱效应（计划 $d = \ph{X.XX}$）；移除关系修饰器应使关系差异性坍塌（计划 \ph{XX}\% 通过率）。
假设是：阶段锚定认知、瞬时情绪和关系状态对临床可解读的动态是共同必要的；上述数字是计划占位符，不是观测。

\paragraph{初始严重度调整下的骨干比较。}
我们登记一个受控比较：在保持架构不变的情况下替换患者/治疗师骨干。
由于各人设起始严重度不同，原始终点比较存在混淆；因此我们报告人设内 z 分数化的终点和一个调整后降幅
\begin{equation}\label{eq:adj}
\mathrm{Adj}\Delta_i \;=\; z_i(\mathrm{T0}) - z_i(\mathrm{S12}),
\end{equation}
其中 $z$ 在比较池内跨人设计算。
Init 列在各人设间差异很大（PHQ-9 上 \ph{X.X}--\ph{X.X}），确认了这一调整的必要性。
表~\ref{tab:backbone} 报告已登记的比较；数值为占位符，待骨干替换批次。

\begin{table}[h]
\centering
\caption{骨干比较（同一架构，替换患者/治疗师骨干）。Init = z 分数化的 T0；Adj$\Delta$ 按式~\ref{eq:adj}。数值为占位符，待已登记批次。}
\label{tab:backbone}
\small
\begin{tabular}{@{}lccc@{}}
\toprule
骨干 & Init (z) & 原始 $\Delta$ & Adj$\Delta$ \\
\midrule
DeepSeek（主线） & \ph{X.XX} & \ph{X.X} & \ph{X.XX} \\
\ph{骨干 B} & \ph{X.XX} & \ph{X.X} & \ph{X.XX} \\
\ph{骨干 C} & \ph{X.XX} & \ph{X.X} & \ph{X.XX} \\
\bottomrule
\end{tabular}
\end{table}

\paragraph{PSI-Bench 读数。}
报告四个语料库（ESConv、AnnoMI、\sys-0922、PatientPSI），每个 8 次会谈 $\times$ 8 条患者回应 $=$ 64 条，共 256 条；PTC 分布与语料库特定参照的一致性达 \ph{XX}\%，情绪分布达 \ph{XX}\%（离线 Qwen3-8B 判断器）。
第五个已下载语料库在两张已报告结果表中均缺席；它是否被评估过，等待归档确认。
逐轮平均 JSD 距离尚未计算，真实人类参照语料库尚未获得——这些是未决事项，不作人类平价声明。

%% ---------------------------------------------------------------
\section{案例研究}\label{app:cases}

\subsection{行为案例}

表~\ref{tab:cases} 把反复出现的行为模式映射到可追溯的证据类别；每个条目链接到已归档的对话与记忆轨迹（\ph{Trace-ID 占位}）。

\begin{table}[h]
\centering
\caption{行为案例类别与证据指针（\ph{Trace-ID} 占位；完整转录存于代码仓库）。}
\label{tab:cases}
\small
\begin{tabular}{@{}lp{10.2cm}@{}}
\toprule
类别 & 模式（证据指针） \\
\midrule
退缩 & 患者拒绝村庄聚会；记忆记下回避（\ph{Trace-ID}）。 \\
作业链 & 提取 $\le$3 条指令；环境反馈；结果记忆在下次会谈被召回（\ph{Trace-ID}）。 \\
微技能门控 & 高防御下选择器切换为确认为先的技能（0--10 量表；\ph{Trace-ID}）。 \\
最小阶段变化 & 变化检测器/阶段规划器/阶段校验器通过后提交单句阶段摘要（\ph{Trace-ID}）。 \\
效价化遭遇 & 负面前缀聊天将痛苦最小化；次日活动下降（\ph{Trace-ID}）。 \\
自然注入 & 未计划的居民问候触发中性前缀（\ph{Trace-ID}）。 \\
记忆阻断会话 & G7 会话运行到阈值；无事件/思想/聊天写入（\ph{Trace-ID}）。 \\
\bottomrule
\end{tabular}
\end{table}

\subsection{长期案例}

\paragraph{案例 1：完全反应轨迹。}
一次 LRN 运行把 PHQ-9 从 \ph{X.X}（T0）移动到 \ph{X.X}（S12），最陡区段在第 \ph{X}--\ph{Y} 次会谈；提交阶段数为 \ph{N}，\ph{XX}\% 通过证据审计（\ph{Trace-ID}）。
\paragraph{案例 2：急性反应保留（持久性已登记）。}
一次 G7 运行在 S12 与 G1 相当（$\Delta$ \ph{X.X}）；其随访读数（FU 处反弹、\ph{XX}\% 严重度重入；主诉图推进停滞于阶段 \ph{标签占位}）已在随访扩展下登记，且尚无同版本随访数据（\ph{Trace-ID}）。
\paragraph{案例 3：社会效价暴露。}
一次 G5 运行恶化 \ph{X.X} 分，退缩模式记忆主导检索；同一人设在 G9 下改善 \ph{X.X}（\ph{Trace-ID}）。

%% ---------------------------------------------------------------
\section{系统提示词}\label{app:prompts}

表~\ref{tab:prompts} 索引各提示词家族；全文见代码仓库，将在数据锁定时以逐字表格附加（\ph{全文占位}）。

\begin{table}[h]
\centering
\caption{提示词家族。全文：\ph{全文占位——代码仓库}。}
\label{tab:prompts}
\small
\begin{tabular}{@{}lp{10.2cm}@{}}
\toprule
ID & 家族 \\
\midrule
P1 & 患者栈渲染器（第 1--4 层 $+$ 记忆与滚动窗口）。 \\
P2 & 变化检测器（严格 JSON；反例约束）。 \\
P3 & 阶段规划器（$\le$100 字符标签；$\le$2 句摘要）。 \\
P4 & 阶段校验器（变化 / 阶段保真度）。 \\
P5 & 医生会谈提示词（12 次会谈大纲）。 \\
P6 & 选择器候选微技能。 \\
P7 & 居民前缀：中性 / 负面 / 积极。 \\
P8 & 量表施测（单轮、无状态）。 \\
P9 & 风险路由器评级。 \\
P10 & 证据审计判断器指令。 \\
\bottomrule
\end{tabular}
\end{table}

%%====================================================================
%%  图（占位）
%%====================================================================

\begin{figure}[t]
\centering
\fbox{\parbox[c][4.2cm][c]{0.93\linewidth}{\centering \emph{\ph{图 1 占位}}\\
(a) 一位患者拒绝一次村庄聚会。\\
(b) 一次医生会谈：第 3.1 次会谈中的苏格拉底式交流。\\
(c) 一次负面效价居民聊天。\\
(d) 下次会谈中的一次作业汇报。}}
\caption{\sys 中的生活。每个场景都描绘了案例研究（表~\ref{tab:cases}）中记录的一个真实行为模式；每个场景的转录轨迹均已归档。}
\label{fig:teaser}
\end{figure}

\begin{figure}[t]
\centering
\fbox{\parbox[c][6.5cm][c]{0.93\linewidth}{\centering \emph{\ph{图 2 占位}}\\
总览：村庄社区 $\rightarrow$ 四层智能体栈 $\rightarrow$ 调度与七臂叠加 $\rightarrow$ 冻结快照与阶段性量表 $\rightarrow$ 审计。}}
\caption{\sys 总览。患者智能体生活在居民村庄中（\S\ref{sec:agent}）；基于规则的调度每 6 步插入一次受控接触（\S\ref{sec:contact}）；七臂叠加仅替换干预（\S\ref{sec:arms}）；冻结快照供给阶段性心理测量与审计（\S\ref{sec:outcome}）。}
\label{fig:overview}
\end{figure}

\begin{figure}[t]
\centering
\fbox{\parbox[c][5.2cm][c]{0.93\linewidth}{\centering \emph{\ph{图 3 占位}}\\
各臂 PHQ-9 / 超短量表轨迹 T0$\rightarrow$S12，跨运行 95\% CI 带。}}
\caption{36 天窗口内的症状轨迹。阴影带表示跨独立运行的 95\% 置信区间；无干预对照与自然接触臂带虚拟会谈标签。数值为占位符。}
\label{fig:traj}
\end{figure}

\begin{figure}[t]
\centering
\fbox{\parbox[c][5.2cm][c]{0.93\linewidth}{\centering \emph{\ph{图 4 占位}}\\
行为--症状相关热图（活动、社交频率、作业完成度 $\times$ 症状变化）。}}
\caption{行为频率与症状变化的相关。作业完成度与社交活动频率分别以 $r = \ph{0.XX}$ 与 $r = \ph{0.XX}$ 与改善相关（\ph{占位}）。}
\label{fig:heatmap}
\end{figure}

\begin{figure}[t]
\centering
\fbox{\parbox[c][5.2cm][c]{0.93\linewidth}{\centering \emph{\ph{图 5 占位}}\\
治疗前后主诉图转移矩阵，G1 对 G7。}}
\caption{治疗前后的主诉图阶段转移。G1 比 G7 更大程度地向残余/恢复阶段推进，而 G7 被阻断的记忆写入损害了持久的认知变化。数值为占位符。}
\label{fig:graph}
\end{figure}

\begin{figure}[t]
\centering
\fbox{\parbox[c][4.6cm][c]{0.93\linewidth}{\centering \emph{\ph{图 6 占位}}\\
按运行阶段的 Token / 调用 / 墙钟时间；每周成本趋势。}}
\caption{计算成本分解。M = 百万 token，K = 千次调用，h = 墙钟小时；数值为占位符，待目标批次。}
\label{fig:cost}
\end{figure}

%%====================================================================
%%  比较表（正文）
%%====================================================================
\begin{table}[t]
\centering
\caption{\sys 与现有方法的比较。$\checkmark$ = 支持，$\times$ = 不支持。真实试验的臂数反映典型实践~\citep{cuijpers2013}。}
\label{tab:comparison}
\footnotesize
\setlength{\tabcolsep}{4pt}
\begin{tabular}{@{}lp{2.1cm}p{1.5cm}p{1.35cm}p{1.0cm}p{1.6cm}@{}}
\toprule
方法 & 媒介 & 时间尺度 & 抑郁表型 & 试验臂 & 标准化结局 \\
\midrule
基于智能体的模型~\citep{bonabeau2002} & 规则智能体 & 数月--数年 & $\times$ & $\times$ & $\times$ \\
生成式智能体~\citep{park2023} & LLM 小镇 & $\sim$2 天 & $\times$ & $\times$ & $\times$ \\
OASIS~\citep{oasis2024} & LLM 社交媒体 & 数天--数周 & $\times$ & $\times$ & $\times$ \\
Agentopia~\citep{agentopia2026} & LLM 小镇 & 10 年 & $\times$ & $\times$ & 仅训练奖励 \\
真实 CBT RCT~\citep{cuijpers2013} & 真实患者 & 数周--数月 & $\checkmark$ & 通常 $\le$4 & $\checkmark$（昂贵） \\
\sys（本文） & LLM 小镇 & 36 仿真日 & $\checkmark$ & 7 $+$ 扩展 & $\checkmark$（PHQ-9 / BDI-II / 短量表） \\
\bottomrule
\end{tabular}
\end{table}

\end{document}
